\section{Source audit and comparison with prior frameworks}

\subsection{Pinned repository state}

The research audit used ProveIt commit
\path{8649990597d73657bde480b6c389de7f85bc0cb0}, committed on
16 November 2026 UTC, which is still 15 November in America/Los\_Angeles at the
time of this work.  The commit itself concerns another project, but includes the
recently placed iPRA framework packets and the new direct derivation library
from preceding commits.

Key source files and roles are:

\begin{longtable}{L{0.48\textwidth}L{0.43\textwidth}}
\toprule
Path & Audited role \\
\midrule
\path{Logic/PrimitiveRecursiveArithmetic/README.md} & Defines iPRA used by the repository and records that proof packets were filed without verification.\\
\path{ConciseIPRA/Lean/ConciseIPRA/NativeClausal.lean} & Defines \(\NFm.\mathsf{QFree}\), \path{unQ}, clauses, and the raw clause calculus.\\
\path{ConciseIPRA/Lean/ConciseIPRA/NativeClausalDiagonal.lean} & Defines \(\KPrf\); proves substitution, cut, soundness to \(\NPrf\), and equivalence with the native kernel.\\
\path{ConciseIPRA/Lean/ConciseIPRA/NativeBiInterpretation.lean} & Defines \path{NFm.toPRA}, \path{NFm.ofPRA}, \path{native_iff_pra}, \path{pra_iff_native}, and round trips.\\
\path{Derivations/README.md} & States the checked status and architecture of direct theorem derivations.\\
\path{Derivations/Lean/IPRADerivations/Core/Basic.lean} & Arithmetic symbols, equations, congruence, \path{ind}, and \path{inst}.\\
\path{Derivations/Lean/IPRADerivations/Core/Tactic.lean} & Universal closure and \path{ipra_ring}.\\
\path{docs/proofs/foundations/ocipra_framework_2026-11-15} & Observation contracts, theorem classes, demand closure, and producer/verifier distinctions.\\
\path{docs/proofs/foundations/ipra_pfig_framework_2026-11-15} & Backward-demand/forward-execution proof graphs and compiler blueprint.\\
\path{docs/proofs/foundations/ipra_proof_carrying_observation_graphs_2026-11-15} & Proof-carrying node architecture, relational paths, and stable observers.\\
\bottomrule
\end{longtable}

\subsection{What is reused}

This report reuses the finite-observation and explicit-resource discipline of
the earlier packets.  In particular, it retains:

\begin{itemize}
\item deterministic, relational, stable-observer, and conditional theorem
classes;
\item backward precision demand and forward replay;
\item bounded trace reflection and relational viability;
\item the separation between producers, verifiers, and finalizers;
\item the warning that cardinality is not a height bound;
\item strict separation of semantic host theorems from object derivations.
\end{itemize}

These ideas are not rebranded as new discoveries.

\subsection{What is new in the present synthesis}

The principal new synthesis is to use the current five-rule \(\KPrf\) clause
calculus as the compiler intermediate representation and to connect it
bidirectionally with the newly checked direct derivation backend.  Earlier
packets generally target \(\Prf\;\PRAAx\) directly or define their own contract
structures.  The present design observes that:

\begin{enumerate}
\item graph composition is exactly \path{KPrf.cut};
\item specialization is exactly \path{KPrf.subst};
\item arbitrary native formulas can be normalized by \path{unQ};
\item export is \path{KPrf.toNPrf} followed by \path{native_iff_pra};
\item direct polynomial derivations can be imported by
\path{pra_iff_native} and \path{kprf_iff_nprf}.
\end{enumerate}

This does not by itself complete the Lean compiler, but it removes most of its
logical burden and gives a sharply bounded implementation project.

\subsection{Environment audit}

The artifact environment contained Python, \path{pdflatex}, \path{latexmk},
\path{pdftoppm}, and PDF inspection tools.  It did not contain \path{lean} or
\path{lake}.  No Lean compilation is therefore claimed.  The package includes
candidate source and a completion plan, and all status files repeat this
limitation.

\subsection{Reproducibility}

From the package root:

\begin{lstlisting}[language={}]
python3 verify.py
./build.sh
\end{lstlisting}

The first command runs standard-library Python tests and checks the manifest.
The second runs tests, regenerates the demo, compiles the PDF, records metadata,
and refreshes checksums.  The build is offline and does not modify the ProveIt
repository.
